\section{Online Adaptation for Continual Tasks}
\label{sec:online_continual_tasks}

The previous sections considered collaborative task planning for a given set of
high-level specifications. In continual operation, this assumption no longer
holds. New tasks may be released while the robots are executing previously
assigned subtasks, and the resulting plan must be updated without discarding
valid progress that has already been achieved. Direct recomputation from the
conjunction of all active formulas is generally undesirable, since the
associated automaton construction and the subsequent multi-robot search may
become prohibitively expensive for long temporal formulas
\cite{baier2008principles,schillinger2018simultaneous,
kantaros2020stylus,luo2021temporal}. For continual tasks, a more suitable
strategy is to update the high-level partial-order structure online and then
re-solve only a finite planning window over the updated structure
\cite{liu2024fast,liu2024time}.

This section develops such an online adaptation mechanism. The first component
is an incremental product of relaxed partially ordered sets, abbreviated as
R-posets, by which a newly released temporal task is merged into the current
global partial order. The second component is a receding-horizon planning
module, inherited from the previous section, that is repeatedly solved on a
finite subset of currently relevant subtasks. The resulting architecture
preserves the correctness of the high-level task logic while maintaining
computational tractability during long-duration execution.

\subsection{Continual task model and planning state}
\label{subsec:continual_model}

Let $t \geq 0$ denote the current planning time. The set of active continual
task formulas is denoted by
$\Phi_t \triangleq \{\varphi_t^1,\cdots,\varphi_t^{m_t}\}$, where each
$\varphi_t^\ell$ is an sc-LTL formula over the atomic propositions introduced
earlier in the chapter. The currently maintained global R-poset is denoted by
\[
\mathcal{P}_t \triangleq (\Omega_t,\preceq_t,\neq_t),
\]
where $\Omega_t$ is the set of active subtasks, $\preceq_t$ is the partial
precedence relation, and $\neq_t$ is the conflict relation. As in the previous
sections, a subtask in $\Omega_t$ may represent either an individual action or
a collaborative action that requires multiple robots.

The online update should account for the fact that part of the current poset
may already have been executed. To this end, let
$\Omega_t^{\mathrm{fin}} \subseteq \Omega_t$ be the set of finished subtasks,
let $\Omega_t^{\mathrm{run}} \subseteq \Omega_t \setminus
\Omega_t^{\mathrm{fin}}$ be the set of subtasks currently under execution, and
let $\Omega_t^{\mathrm{rem}} \triangleq \Omega_t \setminus
\Omega_t^{\mathrm{fin}}$ be the set of remaining subtasks. Moreover, let
$\mathbf{w}_t^{\mathrm{exe}}$ denote the finite execution word that has already
been generated by the system up to time $t$.

\begin{definition}[Continual planning state]
\label{def:continual_state}
The continual planning state at time $t$ is the tuple
\[
\Xi_t \triangleq
\big(\mathcal{P}_t,\Omega_t^{\mathrm{fin}},\Omega_t^{\mathrm{run}},
\mathbf{w}_t^{\mathrm{exe}}\big),
\]
which summarizes the currently maintained task structure, the completed
subtasks, the non-preemptive subtasks already in execution, and the realized
high-level behavior.
\end{definition}

Definition~\ref{def:continual_state} is the minimal information required for
online adaptation. In particular, $\mathbf{w}_t^{\mathrm{exe}}$ records the
executed prefix and therefore distinguishes the already discharged temporal
requirements from the requirements that remain to be satisfied. This distinction
is essential in continual settings, since a newly released task should only
interact with the unfinished portion of the current plan.

\subsection{Online product of R-posets}
\label{subsec:online_poset_product}

Assume that a new continual task is released at time $t$ and is represented by
the sc-LTL formula $\varphi_t^{+}$. Following the same automaton-to-poset
abstraction used earlier in the chapter, let
\[
\mathcal{P}_t^{+} \triangleq
(\Omega_t^{+},\preceq_t^{+},\neq_t^{+})
\]
be an R-poset associated with $\varphi_t^{+}$. The goal is to combine
$\mathcal{P}_t^{+}$ with the current global poset $\mathcal{P}_t$ without
reconstructing the complete task model from scratch.

\begin{definition}[Online product of R-posets]
\label{def:online_poset_product}
Given the continual planning state $\Xi_t$ and the incoming R-poset
$\mathcal{P}_t^{+}$, the online product is the set
\begin{equation}
\mathfrak{P}_t^{+} \triangleq
\mathcal{P}_t \otimes_{\mathbf{w}_t^{\mathrm{exe}}} \mathcal{P}_t^{+},
\label{eq:online_poset_product}
\end{equation}
where each element of $\mathfrak{P}_t^{+}$ is an updated R-poset whose
language is compatible with both the unfinished requirements of
$\mathcal{P}_t$ after the executed prefix $\mathbf{w}_t^{\mathrm{exe}}$ and
the newly released task $\mathcal{P}_t^{+}$.
\end{definition}

The product in \eqref{eq:online_poset_product} is computed over the unfinished
portion of $\mathcal{P}_t$. Consequently, the finished subtasks in
$\Omega_t^{\mathrm{fin}}$ are not reconsidered as decision variables. This is
precisely the mechanism that prevents unnecessary recomputation. In contrast
with direct automaton synchronization, the update is performed only at the
level of the remaining subtasks and their partial relations
\cite{liu2024fast}.

The online product contains two conceptual stages. The first stage constructs a
composed subtask set. Each incoming subtask in $\Omega_t^{+}$ is either merged
with a compatible unfinished subtask in $\Omega_t^{\mathrm{rem}}$ or appended
as a new subtask if no compatible image exists. Compatibility is determined by
the execution label and the self-loop condition inherited from the underlying
automata. Therefore, the composition step does not merely concatenate subtasks;
it also identifies cases in which one future execution can satisfy multiple
specifications simultaneously.

To formalize this point, let
$\mu_t : \Omega_t^{+} \rightarrow \bar{\Omega}_t$ be a composition map, where
$\bar{\Omega}_t$ denotes the composed subtask set produced by the current
update. For each $\omega^{+} \in \Omega_t^{+}$, the image $\mu_t(\omega^{+})$
is either an unfinished subtask in $\Omega_t^{\mathrm{rem}}$ with compatible
labels or a newly created subtask appended to $\bar{\Omega}_t$. The mapping
$\mu_t$ thus records how the new temporal specification is embedded into the
updated global poset.

The second stage refines the relations among subtasks. The inherited precedence
and conflict relations from $\mathcal{P}_t$ and $\mathcal{P}_t^{+}$ are first
lifted to the composed subtask set through $\mu_t$. Additional relations are
then introduced when they are required to preserve self-loop guards or to
remove newly induced conflicts. The updated relation sets are therefore written
as
\begin{equation}
\bar{\preceq}_t \triangleq
\preceq_t \cup \mu_t(\preceq_t^{+}) \cup \preceq_t^{\mathrm{add}},
\qquad
\bar{\neq}_t \triangleq
\neq_t \cup \mu_t(\neq_t^{+}) \cup \neq_t^{\mathrm{add}},
\label{eq:online_relation_update}
\end{equation}
where $\preceq_t^{\mathrm{add}}$ and $\neq_t^{\mathrm{add}}$ denote the
extra precedence and conflict relations introduced during refinement.

Equation~\eqref{eq:online_relation_update} is the key structural update in the
continual planning layer. The first two terms preserve the logic already
encoded in the original posets, while the third term resolves incompatibilities
that appear only after composition. As a result, the updated poset
\[
\bar{\mathcal{P}}_t \triangleq
(\bar{\Omega}_t,\bar{\preceq}_t,\bar{\neq}_t)
\]
remains a valid high-level abstraction of the augmented task specification.
Following \eqref{eq:online_poset_product}, one admissible element of
$\mathfrak{P}_t^{+}$ can be selected as the new global poset
$\mathcal{P}_{t^{+}}$. In practice, the first admissible product is often
sufficient for online reaction, whereas a larger portion of
$\mathfrak{P}_t^{+}$ may be explored when additional computation time is
available \cite{liu2024fast}.

\subsection{Receding-horizon planning over the updated poset}
\label{subsec:online_rh_planning}

Once the global poset has been updated, the next step is to adapt the robot
plans. Replanning over the whole remaining subtask set is unnecessary at every
update and would negate the computational advantage gained in
Subsection~\ref{subsec:online_poset_product}. For this reason, the
receding-horizon planning module introduced in the previous section is used
here as the second layer of online adaptation.

Let $\mathcal{P}_t = (\Omega_t,\preceq_t,\neq_t)$ now denote the updated
global poset after the most recent online product. The set of currently ready
subtasks is defined by
\[
\Omega_t^{\mathrm{ready}} \triangleq
\{\omega \in \Omega_t^{\mathrm{rem}} \setminus \Omega_t^{\mathrm{run}}
\mid \forall \omega' \in \Omega_t,\;
\omega' \preceq_t \omega \Rightarrow
\omega' \in \Omega_t^{\mathrm{fin}}\}.
\]
Hence, a subtask is ready only if all its required predecessors have already
been completed and the subtask is not being executed by another robot team.
This definition excludes temporal violations at the high level and also avoids
duplicate dispatch of a currently running subtask.

Given a horizon length $H \in \mathbb{N}_{>0}$, the planner does not optimize
over all remaining subtasks. Instead, it constructs a finite horizon set
\begin{equation}
\Omega_t^{H} \triangleq
\operatorname{Close}_{H}(\Omega_t^{\mathrm{ready}},\preceq_t),
\label{eq:horizon_set}
\end{equation}
where $\operatorname{Close}_{H}$ returns the first $H$ subtasks obtained by
causal expansion from the ready frontier under the precedence relation
$\preceq_t$.

Equation~\eqref{eq:horizon_set} defines the set of subtasks that are exposed to
the receding-horizon optimizer at the current planning cycle. The subtasks in
$\Omega_t^{H}$ may include both immediately executable subtasks and a small
number of near-future subtasks that are needed to preserve local temporal
consistency. In this way, the planner reasons beyond the current frontier while
avoiding the burden of full-horizon optimization.

The non-preemption requirement is imposed through the currently running set
$\Omega_t^{\mathrm{run}}$. More precisely, let $\mathbb{J}_t^{H}$ denote the
set of feasible horizon plans. This set is defined by
\begin{equation}
\mathbb{J}_t^{H} \triangleq
\left\{
\mathbf{J}\ \middle|\
\begin{array}{l}
\mathbf{J}\ \text{preserves all assignments in } \Omega_t^{\mathrm{run}},\\
\mathbf{J}\ \text{satisfies the relations } \preceq_t \text{ and } \neq_t\\
\text{on the subtasks in } \Omega_t^{H},\\
\mathbf{J}\ \text{satisfies the motion, capability, and}\\
\text{collaborative constraints from the previous section}
\end{array}
\right\},
\label{eq:horizon_feasible_set}
\end{equation}
where $\mathbf{J}$ collects the local plans of all robots over the current
planning horizon.

Equation~\eqref{eq:horizon_feasible_set} shows that the present section does
not alter the underlying receding-horizon solver. Instead, it changes the
high-level input to that solver by updating the poset online and by freezing
the subtasks already in execution. Therefore, the local optimization model from
the previous section can be used without modification, except that it is now
applied to the horizon set $\Omega_t^{H}$ derived from the updated global
poset. After the horizon problem is solved, only the first portion of the
resulting plan is committed for execution, and the rest remains revisable in
later planning cycles.

\subsection{Triggering conditions and state handover}
\label{subsec:trigger_handover}

The preceding subsection describes how replanning is performed after an update.
It remains to specify when replanning is triggered and how information is
passed from one planning cycle to the next. Since continual tasks may arrive at
arbitrary times, the planning process is organized around a sequence of
replanning instants $\tau_0,\tau_1,\cdots$.

Three classes of events are monitored online. The first class is a progress
event. This event is generated when a sufficient fraction of the currently
committed horizon subtasks has been completed. The second class is an arrival
event. This event is generated when a new continual task is released or when an
existing high-level request is modified. The third class is an infeasibility
event. This event is generated when the current horizon plan can no longer be
executed because of execution failure, communication loss, or a newly observed
resource conflict.

Let
\[
r_t \triangleq
\frac{
\left| \Omega_t^{\mathrm{fin}} \cap \Omega_{\tau_\ell}^{H} \right|
}{
\left| \Omega_{\tau_\ell}^{H} \right|
}
\]
be the completion ratio of the horizon plan dispatched at the previous
replanning instant $\tau_\ell$, and let $r_{\mathrm{trig}} \in (0,1]$ be the
prescribed progress threshold. The next replanning instant is then defined as
\begin{equation}
\tau_{\ell+1} \triangleq
\inf \left\{
t > \tau_\ell \,\middle|\,
r_t \geq r_{\mathrm{trig}}
\;\vee\;
\mathsf{E}_t^{\mathrm{new}}
\;\vee\;
\mathsf{E}_t^{\mathrm{fail}}
\right\},
\label{eq:replanning_trigger}
\end{equation}
where $\mathsf{E}_t^{\mathrm{new}}$ denotes the arrival of a new or modified
task, and $\mathsf{E}_t^{\mathrm{fail}}$ denotes infeasibility of the current
horizon plan.

Equation~\eqref{eq:replanning_trigger} unifies periodic progress-driven
replanning with event-driven adaptation. The threshold
$r_{\mathrm{trig}} = 1/2$ is often a reasonable choice in practice, since it
balances reactivity against excessive replanning frequency. More importantly,
the arrival and infeasibility events can trigger replanning immediately, which
is necessary in genuinely dynamic environments.

At each replanning instant, the continual planning state is rolled forward using
the observed execution status of the robots and objects. For each robot
$n \in \mathcal{N}$, let $x_n(\tau_{\ell+1})$ be its measured state at the new
planning instant, and let $J_n^{\mathrm{run}}(\tau_{\ell+1})$ be the frozen
prefix of its local plan that contains the subtask currently under execution,
if any. Then the next planning cycle is initialized from the state snapshot
\[
\Xi_{\tau_{\ell+1}} =
\big(\mathcal{P}_{\tau_{\ell+1}},
\Omega_{\tau_{\ell+1}}^{\mathrm{fin}},
\Omega_{\tau_{\ell+1}}^{\mathrm{run}},
\mathbf{w}_{\tau_{\ell+1}}^{\mathrm{exe}}\big),
\]
together with the measured robot and object states. This state handover keeps
the executed portion of the plan fixed, carries over all non-preemptive
subtasks, and makes the remaining assignments available for revision.

The overall continual adaptation loop can now be summarized as follows.
Whenever \eqref{eq:replanning_trigger} is activated, the incoming task formulas
are first absorbed into the global temporal structure through the online poset
product in \eqref{eq:online_poset_product} and the relation update in
\eqref{eq:online_relation_update}. Next, the ready frontier and finite horizon
set are computed according to \eqref{eq:horizon_set}. Finally, the same
receding-horizon planning module as in the previous section is solved over the
feasible set in \eqref{eq:horizon_feasible_set}, subject to the frozen running
subtasks. In this manner, the robots can incorporate newly released continual
tasks while preserving valid progress and avoiding unnecessary recomputation.
